% !TeX root = part8.tex
\documentclass[../final.tex]{subfiles}

\begin{document}

% При отдельной сборке продолжаем нумерацию полного конспекта.
% Глава о линейной алгебре открывается в part7 и заканчивается прогонкой.
\ifSubfilesClassLoaded{\setcounter{chapter}{6}\setcounter{section}{3}}{}
\captionsetup{font=footnotesize,labelfont=bf,skip=5pt}

% Окончание предыдущей лекции: оценки собственных значений.
% Сам степенной метод в тетради не записан.

\section{Степенной метод}

\subsection{Оценки собственных значений по итерациям}

После установления итерационного процесса отношение норм соседних
приближений даёт оценку собственного значения.

В \rconcept{степенном методе} $\vec x^{(k+1)}=A\vec x^{(k)}$, и
\begin{equation*}
    |\lambda_1|\simeq\frac{\|\vec x^{(k+1)}\|}{\|\vec x^{(k)}\|},
\end{equation*}
где $\lambda_1$ --- наибольшее по модулю собственное значение.

В \rconcept{обратном степенном методе} $\vec x^{(k+1)}=A^{-1}\vec x^{(k)}$.
Наибольшее по модулю собственное значение матрицы $A^{-1}$ равно
$1/\lambda_n$, где $\lambda_n$ --- наименьшее по модулю
собственное значение $A$:
\begin{equation*}
    \frac{1}{|\lambda_n|}\simeq
    \frac{\|\vec x^{(k+1)}\|}{\|\vec x^{(k)}\|}.
\end{equation*}

В \rconcept{методе обратных итераций со сдвигом}
$\vec x^{(k+1)}=\bigl(A-\lambda_i^{(0)}E\bigr)^{-1}\vec x^{(k)}$,
где $\lambda_i^{(0)}$ --- начальное приближение к $\lambda_i$.
Процесс сходится к собственному значению, ближайшему к сдвигу:
\begin{equation*}
    \frac{1}{\bigl|\lambda_i-\lambda_i^{(0)}\bigr|}\simeq
    \frac{\|\vec x^{(k+1)}\|}{\|\vec x^{(k)}\|}
    \quad\Longrightarrow\quad
    \lambda_i\simeq\lambda_i^{(0)}\pm
    \frac{\|\vec x^{(k)}\|}{\|\vec x^{(k+1)}\|}.
\end{equation*}
Отношение норм определяет только модуль. \rnote{Знак поправки
определяется тем, сохраняет ли $\vec x^{(k+1)}$ направление
$\vec x^{(k)}$ или меняет его на противоположное.}

% ============================================================
% Лекция 27.10.2026
% Решение задач линейной алгебры: поиск собственных значений
% и собственных векторов
% ============================================================

\section{Метод вращений Якоби}

\subsection{Идея метода}

Пусть $A$ --- \rassumption{вещественная симметричная матрица}.
Тогда существует \rconcept{ортогональная матрица} $O$, $O^TO=E$,
приводящая $A$ к диагональному виду:
\begin{equation*}
    \Lambda=O^TAO=\operatorname{diag}(\lambda_1,\ldots,\lambda_n).
\end{equation*}
На диагонали $\Lambda$ стоят собственные значения $A$,
а столбцы $O$ --- соответствующие собственные векторы.

Метод Якоби строит $O$ итерационно, как произведение элементарных
вращений. Каждое вращение обнуляет один внедиагональный элемент.

\subsection{Шаг итерации}

Пусть $A^{(k)}$ --- приближение к диагональной матрице $\Lambda$
на $k$-й итерации, $A^{(0)}=A$. Выберем наибольший по модулю
внедиагональный элемент:
\begin{equation*}
    \bigl|a^{(k)}_{i_kj_k}\bigr|=\max_{i<j}\bigl|a^{(k)}_{ij}\bigr|,
\end{equation*}
где $i_k$ и $j_k$ --- номер строки и номер столбца максимального
элемента на $k$-й итерации. Благодаря симметрии достаточно
просматривать элементы над диагональю.

Следующее приближение получаем преобразованием
\begin{equation*}
    A^{(k+1)}=V_{i_kj_k}^TA^{(k)}V_{i_kj_k},
\end{equation*}
где $V_{i_kj_k}$ --- \rconcept{матрица вращения} в плоскости
координат $i_k$, $j_k$ на угол $\theta_k$. Она совпадает с единичной
матрицей всюду, кроме четырёх элементов на пересечениях строк
и столбцов с номерами $i_k$ и $j_k$:
\begin{equation*}
    V_{i_kj_k}=
    \begin{pmatrix}
        1      &        &               &        &               &        &   \\
               & \ddots &               &        &               &        &   \\
               &        & \cos\theta_k  & \cdots & -\sin\theta_k &        &   \\
               &        & \vdots        & \ddots & \vdots        &        &   \\
               &        & \sin\theta_k  & \cdots & \cos\theta_k  &        &   \\
               &        &               &        &               & \ddots &   \\
               &        &               &        &               &        & 1
    \end{pmatrix}
    \begin{matrix}
        \\ \\ \leftarrow i_k \\ \\ \leftarrow j_k \\ \\ \\
    \end{matrix}
\end{equation*}
Матрица $V_{i_kj_k}$ ортогональна, поэтому преобразование
сохраняет симметрию и собственные значения.

\subsection{Выбор угла поворота}

Угол $\theta_k$ выбираем так, чтобы элемент на месте
$(i_k,j_k)$ обратился в нуль и матрица $A^{(k+1)}$ стала
ближе к диагональной. Непосредственное перемножение даёт
\begin{equation*}
    a^{(k+1)}_{i_kj_k}
    =\bigl(\cos^2\theta_k-\sin^2\theta_k\bigr)a^{(k)}_{i_kj_k}
    +\cos\theta_k\sin\theta_k
    \bigl(a^{(k)}_{j_kj_k}-a^{(k)}_{i_ki_k}\bigr)=0.
\end{equation*}
Так как $\cos^2\theta-\sin^2\theta=\cos2\theta$ и
$2\cos\theta\sin\theta=\sin2\theta$, отсюда
\begin{equation*}
    \tg2\theta_k=
    \frac{2a^{(k)}_{i_kj_k}}{a^{(k)}_{i_ki_k}-a^{(k)}_{j_kj_k}}
    \quad\Longrightarrow\quad
    \theta_k=\frac12\arctg
    \frac{2a^{(k)}_{i_kj_k}}{a^{(k)}_{i_ki_k}-a^{(k)}_{j_kj_k}}.
\end{equation*}
При $a^{(k)}_{i_ki_k}=a^{(k)}_{j_kj_k}$ берут $\theta_k=\pm\pi/4$
со знаком $a^{(k)}_{i_kj_k}$.

\remark
Обнулённый элемент на следующих шагах может снова стать ненулевым.
Однако сумма квадратов внедиагональных элементов при каждом вращении
уменьшается ровно на $2\bigl(a^{(k)}_{i_kj_k}\bigr)^2$, поэтому
\rresult{матрицы $A^{(k)}$ сходятся к диагональной}.

\subsection{Собственные значения и векторы}

Раскроем последовательность преобразований:
\begin{align*}
    A^{(k+1)}
    &=V_{i_kj_k}^TA^{(k)}V_{i_kj_k}
    =V_{i_kj_k}^TV_{i_{k-1}j_{k-1}}^TA^{(k-1)}
    V_{i_{k-1}j_{k-1}}V_{i_kj_k}=\ldots\\
    &=\underbrace{V_{i_kj_k}^T\cdots V_{i_1j_1}^T}_{\simeq\,O^T}\,
    A\,
    \underbrace{V_{i_1j_1}\cdots V_{i_kj_k}}_{\simeq\,O}.
\end{align*}
Итерации продолжают, пока внедиагональные элементы не станут
меньше заданной точности $\varepsilon$. Тогда $A^{(k+1)}$
\rnote{практически диагональна}, и
\begin{itemize}
    \item собственные значения --- это \rresult{диагональные элементы}
    $A^{(k+1)}$;
    \item собственные векторы --- это \rresult{столбцы матрицы}
    $V_{i_1j_1}\cdots V_{i_kj_k}\simeq O$.
\end{itemize}

% ============================================================
% Решение уравнения теплопроводности
% ============================================================

\chapter{Решение уравнения теплопроводности}
\rsubtitle{Явная, неявная схемы и схема Кранка--Николсон}

\section{Разностная схема и её свойства}

\subsection{Невязка и порядок аппроксимации}

Пусть $L_hy=f$ --- \rconcept{численная схема} решения краевой задачи
для функции $y(x,t)$ с граничными условиями
\begin{equation*}
    \begin{cases}
        y(0,t)=g_0(t),\\
        y(L,t)=g_1(t)
    \end{cases}
\end{equation*}
и начальным условием $y(x,0)=u_0(x)$.
Здесь $h$ и $\tau$ --- шаги сетки по пространству и по времени.

\rconcept{Невязкой} $\psi_{h,\tau}$ численной аппроксимации
дифференциального уравнения называют результат подстановки
точного решения $u$ дифференциального уравнения в численную
аппроксимацию:
\begin{equation*}
    \psi_{h,\tau}=L_hu-f.
\end{equation*}

Говорят, что численная схема \rconcept{аппроксимирует} дифференциальное
уравнение с порядком $p$ по пространству и $s$ по времени, если
\begin{equation*}
    \psi_{h,\tau}=O(h^p+\tau^s).
\end{equation*}

\subsection{Устойчивость}

Разностная схема называется \rconcept{устойчивой}, если существуют
$h_0$, $\tau_0$ и $\delta$ такие, что при всех $h<h_0$, $\tau<\tau_0$
и любом возмущении $\|\varepsilon\|<\delta$ решение возмущённой задачи
\begin{equation*}
    L_h\tilde y=f+\varepsilon
\end{equation*}
удовлетворяет оценке
\begin{equation*}
    \|\tilde y-y\|<C\|\varepsilon\|.
\end{equation*}
Малое возмущение правой части приводит к малому изменению решения.

Если схема устойчива при любых $h$ и $\tau$
($h_0,\tau_0\to\infty$), то она \rconcept{абсолютно устойчива}.
Если схема устойчива только при определённых соотношениях
$h$ и $\tau$, то она \rconcept{условно устойчива}.

\subsection{Сходимость}

\textbf{Теорема.} Если разностная схема аппроксимирует дифференциальное
уравнение с порядком точности $p$ по пространству и $s$ по времени
и схема устойчива, то имеет место оценка
\begin{equation*}
    \|u-y\|<C_1h^p+C_2\tau^s.
\end{equation*}
Таким образом, \rresult{аппроксимация и устойчивость вместе
обеспечивают сходимость} численного решения к точному.

\section{Уравнение теплопроводности}

\subsection{Постановка задачи}

\rconcept{Уравнение теплопроводности} (диффузии) имеет вид
\begin{equation*}
    \frac{\partial u}{\partial t}
    =a\frac{\partial^2u}{\partial x^2}+f(x,t,u),
\end{equation*}
где $a$ --- коэффициент диффузии (температуропроводности).
Функция $u$ --- это температура $T$ в уравнении теплопроводности
или концентрация $n$ в уравнении диффузии.

Решение ищем в области
\begin{equation*}
    0\le x\le 1,\qquad 0\le t\le T.
\end{equation*}
Отрезок $[0,L]$ приводится к $[0,1]$ \rassumption{заменой переменных}.

Задаётся \rconcept{начальное условие}
\begin{equation*}
    u(x,0)=u_0(x)
\end{equation*}
и граничные условия на концах отрезка.
\rconcept{Граничные условия I рода} (Дирихле) задают значения функции:
\begin{equation*}
    \begin{cases}
        u(0,t)=g_0(t),\\
        u(1,t)=g_1(t).
    \end{cases}
\end{equation*}
\rconcept{Граничные условия II рода} (Неймана) задают значения
производной, то есть поток:
\begin{equation*}
    \begin{cases}
        \dfrac{\partial u}{\partial x}(0,t)=g_0(t),\\[0.8em]
        \dfrac{\partial u}{\partial x}(1,t)=g_1(t).
    \end{cases}
\end{equation*}

\subsection{Переход к однородным граничным условиям}

От неоднородных граничных условий можно перейти к однородным,
выделив из решения известную функцию, которая сама удовлетворяет
граничным условиям.

\subsubsection{Условия I рода}

Положим
\begin{equation*}
    u(x,t)=v(x,t)+(1-x)g_0(t)+xg_1(t).
\end{equation*}
Тогда
\begin{align*}
    x=0:&\quad u(0,t)=v(0,t)+g_0(t)=g_0(t)
    \;\Longrightarrow\; v(0,t)=0,\\
    x=1:&\quad u(1,t)=v(1,t)+g_1(t)=g_1(t)
    \;\Longrightarrow\; v(1,t)=0.
\end{align*}

\subsubsection{Условия II рода}

Ищем добавку в виде квадратичного полинома по $x$:
\begin{equation*}
    u(x,t)=v(x,t)+\alpha\frac{x^2}{2}+\beta x.
\end{equation*}
Потребуем $v'_x=0$ на обеих границах:
\begin{equation*}
    \begin{cases}
        u'_x\big|_{x=0}=v'_x\big|_{x=0}+\alpha\cdot0+\beta=g_0(t),\\
        u'_x\big|_{x=1}=v'_x\big|_{x=1}+\alpha+\beta=g_1(t),
    \end{cases}
    \quad\Longrightarrow\quad
    \alpha=g_1-g_0,\qquad\beta=g_0.
\end{equation*}
Следовательно,
\begin{equation*}
    u(x,t)=v(x,t)+\frac{x^2}{2}\bigl[g_1(t)-g_0(t)\bigr]+xg_0(t).
\end{equation*}

\remark
Добавка зависит от $t$ и, для условий II рода, не является линейной
по $x$. Поэтому функция $v$ удовлетворяет \rnote{тому же уравнению
с изменённой правой частью}:
\begin{equation*}
    \frac{\partial v}{\partial t}=a\frac{\partial^2v}{\partial x^2}
    +\tilde f,\qquad
    \tilde f=f-\frac{\partial w}{\partial t}
    +a\frac{\partial^2w}{\partial x^2},
\end{equation*}
где $w=u-v$ --- выделенная добавка. Начальное условие также
пересчитывается: $v(x,0)=u_0(x)-w(x,0)$.

\section{Конечно-разностная аппроксимация}

\subsection{Сетка}

В области $0\le x\le L$, $0\le t\le T$ введём
\rconcept{равномерную сетку}:
\begin{equation*}
    x_j=jh,\quad h=\frac LN,\quad 0\le j\le N;
    \qquad
    t_n=n\tau,\quad \tau=\frac TM,\quad 0\le n\le M.
\end{equation*}
Значение численного решения в узле обозначим
$y_j^n\simeq u(x_j,t_n)$: \rnote{номер шага по времени пишется сверху,
номер узла по пространству --- снизу}.

\begin{rbox}[breakable=false,left=18pt,right=18pt,top=14pt,bottom=10pt]{[]}
    \centering
    \begin{tikzpicture}[x=1.05cm,y=0.75cm,
        every node/.style={text=rtext,font=\small}]
        \draw[rtext!30] (0,0) grid (8,5);
        \draw[rtext,-{Stealth[length=2.2mm]}] (0,0)--(8.8,0) node[right] {$x$};
        \draw[rtext,-{Stealth[length=2.2mm]}] (0,0)--(0,5.7) node[left] {$t$};
        \draw[p3,line width=1.3pt] (0,0)--(8,0);
        \draw[red3,line width=1.3pt] (0,0)--(0,5);
        \draw[red3,line width=1.3pt] (8,0)--(8,5);
        \foreach \i in {0,...,8}{
            \foreach \n in {0,...,5}{
                \fill[rtext] (\i,\n) circle(1.3pt);
            }
        }
        \node[below=3pt] at (0,0) {$0$};
        \node[below=3pt] at (4,0) {$x_j$};
        \node[below=3pt] at (8,0) {$L$};
        \node[left=3pt] at (0,3) {$t_n$};
        \node[left=3pt] at (0,5) {$T$};
        \draw[rtext,{Stealth[length=1.6mm]}-{Stealth[length=1.6mm]}]
            (5,-0.45)--(6,-0.45) node[midway,below] {$h$};
        \draw[rtext,{Stealth[length=1.6mm]}-{Stealth[length=1.6mm]}]
            (8.45,3)--(8.45,4) node[midway,right] {$\tau$};
        \fill[rc4] (4,3) circle(2.2pt);
        \node[rc4,above right] at (4,3) {$y_j^n$};
        \node[p3,above] at (2,0) {$u_0(x)$};
        \node[red3,left,align=right] at (-0.1,1.5) {$g_0(t)$};
        \node[red3,right] at (8.1,1.5) {$g_1(t)$};
    \end{tikzpicture}
    \captionof{figure}{Равномерная сетка: на нижнем слое задано начальное
    условие, на боковых границах --- граничные условия.}
    \label{fig:ovf8_grid}
\end{rbox}

\subsection{Явная схема}

Заменим производные разностями:
\begin{equation*}
    \frac{\partial u}{\partial t}\to\frac{y_j^{n+1}-y_j^n}{\tau},
    \qquad
    \frac{\partial^2u}{\partial x^2}\to
    \frac{y_{j+1}^n-2y_j^n+y_{j-1}^n}{h^2}.
\end{equation*}
Пусть \rassumption{$a=1$, $f=0$}, то есть
\begin{equation*}
    \frac{\partial u}{\partial t}=\frac{\partial^2u}{\partial x^2}.
\end{equation*}
Вторую производную по $x$ можно взять как на $n$-м, так и на
$(n+1)$-м временном слое. Выбор $n$-го слоя даёт
\rconcept{явную схему}:
\begin{equation*}
    \frac{y_j^{n+1}-y_j^n}{\tau}
    =\frac{y_{j+1}^n-2y_j^n+y_{j-1}^n}{h^2}.
\end{equation*}
Значение $y_j^{n+1}$ на новом слое выражается явно через три
известных значения на предыдущем слое.

\subsection{Неявная схема}

Если вторую производную по $x$ взять на $(n+1)$-м слое, получаем
\rconcept{неявную схему}:
\begin{equation*}
    \frac{y_j^{n+1}-y_j^n}{\tau}
    =\frac{y_{j+1}^{n+1}-2y_j^{n+1}+y_{j-1}^{n+1}}{h^2}.
\end{equation*}
Каждое уравнение связывает три неизвестных значения нового слоя,
поэтому на каждом шаге по времени решается система
с трёхдиагональной матрицей, например методом прогонки.

\subsection{Схема Кранка--Николсон}

Полусумма явной и неявной аппроксимаций второй производной даёт
\rconcept{схему Кранка--Николсон}:
\begin{equation*}
    \frac{y_j^{n+1}-y_j^n}{\tau}
    =\frac12\,\frac{y_{j+1}^n-2y_j^n+y_{j-1}^n}{h^2}
    +\frac12\,\frac{y_{j+1}^{n+1}-2y_j^{n+1}+y_{j-1}^{n+1}}{h^2}.
\end{equation*}
Эта схема имеет \rresult{порядок аппроксимации $O(\tau^2+h^2)$}.

\subsection{Шаблоны схем}

Набор узлов, которые входят в одно разностное уравнение, называется
\rconcept{шаблоном схемы}.

\begin{rbox}[breakable=false,left=18pt,right=18pt,top=14pt,bottom=10pt]{[]}
    \centering
    \begin{tikzpicture}[x=0.95cm,y=1.05cm,
        every node/.style={text=rtext,font=\small}]
        \foreach \dx/\name in {0/{Явная},5/{Неявная},10/{Кранк--Николсон}}{
            \begin{scope}[xshift=\dx cm]
                \draw[rtext!30] (-1.5,0)--(1.5,0);
                \draw[rtext!30] (-1.5,1)--(1.5,1);
                \node[left] at (-1.5,0) {$n$};
                \node[left] at (-1.5,1) {$n+1$};
                \foreach \xx/\lab in {-1/{j-1},0/{j},1/{j+1}}{
                    \draw[rtext!30,dashed] (\xx,-0.15)--(\xx,1.15);
                    \node[below=4pt] at (\xx,-0.1) {$\lab$};
                }
                \node[above=8pt] at (0,1.15) {\name};
            \end{scope}
        }
        % Явная схема.
        \draw[p3,line width=1.1pt] (-1,0)--(1,0);
        \draw[red3,line width=1.1pt] (0,0)--(0,1);
        \foreach \xx in {-1,0,1}{\fill[p3] (\xx,0) circle(2.2pt);}
        \fill[red3] (0,1) circle(2.2pt);
        % Неявная схема.
        \begin{scope}[xshift=5cm]
            \draw[red3,line width=1.1pt] (-1,1)--(1,1);
            \draw[p3,line width=1.1pt] (0,0)--(0,1);
            \foreach \xx in {-1,0,1}{\fill[red3] (\xx,1) circle(2.2pt);}
            \fill[p3] (0,0) circle(2.2pt);
        \end{scope}
        % Схема Кранка--Николсон.
        \begin{scope}[xshift=10cm]
            \draw[p3,line width=1.1pt] (-1,0)--(1,0);
            \draw[red3,line width=1.1pt] (-1,1)--(1,1);
            \draw[rtext,line width=0.8pt] (0,0)--(0,1);
            \foreach \xx in {-1,0,1}{
                \fill[p3] (\xx,0) circle(2.2pt);
                \fill[red3] (\xx,1) circle(2.2pt);
            }
        \end{scope}
    \end{tikzpicture}
    \captionof{figure}{Шаблоны схем: голубые узлы лежат на известном
    слое $n$, розовые --- на новом слое $n+1$.}
    \label{fig:ovf8_stencils}
\end{rbox}

\section{Порядок аппроксимации}

Рассмотрим уравнение с \rassumption{$a=1$, $f=0$}.
Подставим в разностные схемы точное решение $u$ и разложим
все значения по формуле Тейлора в узле $(x_j,t_n)$.
Производные без верхнего индекса берутся в этом узле.

\subsection{Явная схема}

Невязка явной схемы равна
\begin{equation*}
    \psi_{h,\tau}=L_hu-f
    =\frac{u_j^{n+1}-u_j^n}{\tau}
    -\frac{u_{j+1}^n-2u_j^n+u_{j-1}^n}{h^2}.
\end{equation*}
Разложения по времени и по пространству:
\begin{align*}
    u_j^{n+1}&=u_j^n+\frac{\partial u}{\partial t}\tau
    +\frac{\partial^2u}{\partial t^2}\frac{\tau^2}{2}+O(\tau^3),\\
    u_{j\pm1}^n&=u_j^n\pm\frac{\partial u}{\partial x}h
    +\frac{\partial^2u}{\partial x^2}\frac{h^2}{2}
    \pm\frac{\partial^3u}{\partial x^3}\frac{h^3}{3!}
    +\frac{\partial^4u}{\partial x^4}\frac{h^4}{4!}+O(h^5).
\end{align*}
При сложении $u_{j+1}^n$ и $u_{j-1}^n$ нечётные члены сокращаются:
\begin{equation*}
    \frac{u_{j+1}^n-2u_j^n+u_{j-1}^n}{h^2}
    =\frac{\partial^2u}{\partial x^2}
    +\frac{h^2}{12}\frac{\partial^4u}{\partial x^4}+O(h^4).
\end{equation*}
Подставляя, получаем
\begin{align*}
    \psi_{h,\tau}
    &=\frac{\partial u}{\partial t}
    +\frac{\partial^2u}{\partial t^2}\frac{\tau}{2}
    -\frac{\partial^2u}{\partial x^2}
    -\frac{h^2}{12}\frac{\partial^4u}{\partial x^4}
    +O(\tau^2+h^4)\\
    &=\frac{\tau}{2}\frac{\partial^2u}{\partial t^2}
    -\frac{h^2}{12}\frac{\partial^4u}{\partial x^4}
    +O(\tau^2+h^4)=O(\tau+h^2).
\end{align*}
Здесь $\dfrac{\partial u}{\partial t}-\dfrac{\partial^2u}{\partial x^2}=0$,
так как $u$ --- решение дифференциального уравнения.
Явная схема имеет \rresult{первый порядок точности по $t$
и второй порядок по $x$}.

\subsection{Неявная схема}

Для неявной схемы пространственная разность берётся на слое $n+1$:
\begin{equation*}
    \psi_{h,\tau}
    =\frac{\partial u}{\partial t}
    +\frac{\partial^2u}{\partial t^2}\frac{\tau}{2}
    -\left[\frac{\partial^2u^{n+1}}{\partial x^2}
    +\frac{h^2}{12}\frac{\partial^4u^{n+1}}{\partial x^4}\right]
    +O(\tau^2+\tau h^2+h^4).
\end{equation*}
Перенесём вторую производную с нового слоя в узел $(x_j,t_n)$:
\begin{equation*}
    \frac{\partial^2u^{n+1}}{\partial x^2}
    =\frac{\partial^2}{\partial x^2}
    \left(u+\tau\frac{\partial u}{\partial t}\right)+O(\tau^2)
    =\frac{\partial^2u}{\partial x^2}
    +\tau\frac{\partial}{\partial t}\frac{\partial^2u}{\partial x^2}
    +O(\tau^2)
    =\frac{\partial^2u}{\partial x^2}
    +\tau\frac{\partial^2u}{\partial t^2}+O(\tau^2).
\end{equation*}
В последнем переходе использовано уравнение
$\dfrac{\partial^2u}{\partial x^2}=\dfrac{\partial u}{\partial t}$.
Тогда
\begin{align*}
    \psi_{h,\tau}
    &=\frac{\partial u}{\partial t}
    +\frac{\tau}{2}\frac{\partial^2u}{\partial t^2}
    -\frac{\partial^2u}{\partial x^2}
    -\tau\frac{\partial^2u}{\partial t^2}
    -\frac{h^2}{12}\frac{\partial^4u}{\partial x^4}
    +O(\tau^2+\tau h^2+h^4)\\
    &=-\frac{\tau}{2}\frac{\partial^2u}{\partial t^2}
    -\frac{h^2}{12}\frac{\partial^4u}{\partial x^4}
    +O(\tau^2+\tau h^2+h^4)=O(\tau+h^2).
\end{align*}
Неявная схема также имеет \rresult{первый порядок точности по $t$
и второй порядок по $x$}.

\subsection{Схема Кранка--Николсон}

Невязка схемы Кранка--Николсон равна полусумме невязок явной
и неявной схем. Главные члены по времени у них равны
$\pm\dfrac{\tau}{2}\dfrac{\partial^2u}{\partial t^2}$
и при усреднении сокращаются:
\begin{equation*}
    \psi_{h,\tau}
    =-\frac{h^2}{12}\frac{\partial^4u}{\partial x^4}
    +O(\tau^2+\tau h^2+h^4)=O(\tau^2+h^2).
\end{equation*}
Отсюда \rresult{второй порядок точности по $t$ и по $x$}.

\end{document}
